\documentclass[11pt,twocolumn]{article}
\usepackage{times}
\usepackage{latexsym}
\usepackage{microtype}
\usepackage{booktabs}
\usepackage{multirow}
\usepackage{graphicx}
\usepackage{amsmath}
\usepackage{url}
\usepackage{array}
\usepackage[margin=1in]{geometry}
\usepackage[numbers,sort&compress]{natbib}
\usepackage{hyperref}
\usepackage{setspace}
\usepackage{titlesec}

\titlespacing\section{0pt}{6pt plus 2pt minus 1pt}{3pt plus 1pt minus 1pt}
\titlespacing\subsection{0pt}{4pt plus 2pt minus 1pt}{2pt plus 1pt minus 1pt}

\title{\large\textbf{Monolingual vs.\ Multilingual Fine-Tuning of AfroXLMR for\\
Sentiment Analysis Across Five African Languages:\\
A Code-Mixing Diagnostic Study}}

\author{
  \small Praises Obi\textsuperscript{1} \quad Ishe Allen Chihobo\textsuperscript{1} \\
  \small \textsuperscript{1}Department of Computer Science, University of Pretoria | Group 60 \\
  \small \texttt{\{u26819661, u22592238\}@tuks.co.za}
}

\date{}

\begin{document}
\maketitle

\begin{abstract}
\small
Sentiment analysis for African languages is constrained by limited annotated data, structural
diversity, and pervasive English code-mixing. We compare four systems on five AfriSenti languages
(Hausa, Yor\`{u}b\'{a}, Igbo, Nigerian Pidgin, Swahili): TF-IDF logistic regression, mBERT, and
two fine-tuning regimes of AfroXLMR-base---Language-Adaptive Fine-Tuning (LAFT) and Multilingual
Adaptive Fine-Tuning (MAFT). Our novel contribution is a token-level code-mixing index (CMI)
filter partitioning each test set into monolingual and code-mixed subsets. LAFT achieves the
highest weighted F1 on resource-rich Hausa (76.34); MAFT outperforms LAFT on low-resource Swahili
(62.51 vs.\ 60.86), confirming a resource-driven crossover. Code-mixing effects are
language-dependent, with Yor\`{u}b\'{a} improving on mixed tweets (+6.13 F1) while Hausa degrades
modestly.
\end{abstract}

\section{Introduction}

Africa's 2,000+ languages remain severely underrepresented in NLP systems
\cite{adebara-abdulmageed-2022-towards}. Sentiment analysis---the automatic classification of
opinion as positive, negative, or neutral---is among the highest-demand downstream tasks, and
its failure on African languages translates into absent or biased insights for over one billion
speakers in health, political, and commercial domains.

The AfriSenti corpus \cite{muhammad-etal-2023-afrisenti} provides annotated sentiment data for
14 African languages, and AfroXLMR \cite{alabi-etal-2022-adapting} extends XLM-R with
African-language pretraining. Yet two questions remain open: (1) whether monolingual (LAFT) or
joint multilingual (MAFT) fine-tuning of AfroXLMR is preferable, and whether this depends on
resource level; and (2) whether code-mixing with English systematically degrades performance.

Our contributions are: (1) a controlled four-way comparison of TF-IDF, mBERT, LAFT, and MAFT
across five typologically diverse languages; (2) a token-level CMI filter producing separate
evaluations on monolingual and code-mixed test subsets; (3) empirical confirmation of a
resource-driven LAFT/MAFT crossover; and (4) evidence that code-mixing effects are
language-dependent. This work directly addresses African language underrepresentation and provides
practitioners with evidence-based model-selection guidance.

\section{Background}

\citet{muhammad-etal-2022-naijasenti} introduced NaijaSenti, establishing sentiment annotation
protocols and documenting pervasive code-mixing in Nigerian social media. They showed that standard
multilingual models struggle particularly with Nigerian Pidgin due to its English creole structure.
\citet{muhammad-etal-2023-afrisenti} extended this to AfriSenti (14 languages), and the
AfriSenti-SemEval shared task \cite{muhammad-etal-2023-semeval} benchmarked 53 systems, with top
results between 65--80 weighted F1, but without disaggregating by code-mixing rate.

\citet{mabokela-schlippe-2022} and \citet{mokhosi-etal-2024} demonstrated that models trained on
high-resource languages transfer poorly to morphologically rich Bantu and Sotho languages,
motivating language-specific training strategies. \citet{devlin-etal-2019-bert} introduced mBERT
across 104 languages, but its African-language training data is negligible relative to European
languages. \citet{alabi-etal-2022-adapting} introduced AfroXLMR, adapted via continued pretraining
on 17 African languages, achieving state-of-the-art results on multiple African benchmarks.

\textbf{Gap.} No study has directly compared LAFT and MAFT of the same AfroXLMR model across
multiple language families on AfriSenti, nor applied a principled code-mixing diagnostic to
evaluate performance degradation attributable to English code-mixing. We address both gaps.

\section{Methodology}

\subsection{Dataset}

We use AfriSenti \cite{muhammad-etal-2023-afrisenti}, publicly available under Creative
Commons licence. Five languages were selected to represent distinct typological families and
resource levels (Table~\ref{tab:dataset}).

\begin{table}[t]
\centering
\small
\setlength{\tabcolsep}{3pt}
\begin{tabular}{lrrrr}
\toprule
\textbf{Language} & \textbf{Train} & \textbf{Val} & \textbf{Test} & \textbf{Imb.} \\
\midrule
Hausa (hau, Afro-As.)   & 14,172 & 2,677 & 5,303 & 1.07$\times$ \\
Yor\`{u}b\'{a} (yor, NC) &  8,522 & 2,090 & 4,515 & 1.89$\times$ \\
Igbo (ibo, NC)           & 10,192 & 1,841 & 3,682 & 1.73$\times$ \\
Nig. Pidgin (pcm, Cre.)  &  5,121 & 1,281 & 4,154 & 45.0$\times$ \\
Swahili (swa, NC)        &  1,810 &   453 &   748 & 5.61$\times$ \\
\bottomrule
\end{tabular}
\caption{Dataset statistics. Imbalance = max/min class count in training split. NC = Niger-Congo, Afro-As. = Afro-Asiatic, Cre. = English Creole.}
\label{tab:dataset}
\end{table}

Labels are positive, neutral, or negative. Nigerian Pidgin's extreme 45$\times$ imbalance
(72 neutral vs.\ 3,241 negative training examples) and Swahili's small training set (1,810)
represent the primary data challenges.

\subsection{Code-Mixing Diagnostic Filter}

For each test tweet, we compute an English Token Ratio (ETR):
\[
\text{ETR}(t) = \frac{|\text{English tokens}(t)|}{|\text{valid tokens}(t)|}
\]
where valid tokens exclude @mentions, URLs, numerals, and tokens $<$3 characters.
Per-token language detection uses \texttt{langdetect}. A tweet is code-mixed if ETR $> 0.20$.
This threshold follows prior code-mixing literature and was not tuned on held-out data.
Table~\ref{tab:cmi} shows the resulting distribution.

\begin{table}[t]
\centering
\small
\setlength{\tabcolsep}{3pt}
\begin{tabular}{lrrr}
\toprule
\textbf{Language} & \textbf{Test $n$} & \textbf{Mono.} & \textbf{Mixed (\%)} \\
\midrule
Hausa           & 5,303 & 5,174 & 129 (2.4\%)  \\
Yor\`{u}b\'{a} & 4,515 & 4,352 & 163 (3.6\%)  \\
Igbo            & 3,682 & 3,450 & 232 (6.3\%)  \\
Nigerian Pidgin & 4,154 & 2,533 & 1,621 (39.0\%) \\
Swahili         &   748 &   736 &  12 (1.6\%)\textsuperscript{*}  \\
\bottomrule
\end{tabular}
\caption{Code-mixing diagnostic split. \textsuperscript{*}Swahili mixed subset ($n$=12) is
too small for robust evaluation.}
\label{tab:cmi}
\end{table}

\subsection{Systems}

\textbf{TF-IDF + LR.} Unigram/bigram TF-IDF (50k features, sublinear TF, diacritics preserved)
with logistic regression (L2, $C$=1.0, balanced class weights). Non-neural lower bound.

\textbf{mBERT.} \texttt{bert-base-multilingual-cased} \cite{devlin-etal-2019-bert} fine-tuned
per language. Standard cross-lingual reference.

\textbf{LAFT.} \texttt{Davlan/afro-xlmr-base} \cite{alabi-etal-2022-adapting} fine-tuned
independently on each language's training set. Monolingual specialisation.

\textbf{MAFT.} AfroXLMR-base fine-tuned on the concatenation of all five training sets
(39,817 examples), then evaluated per language. Tests cross-lingual parameter sharing.

\subsection{Training Setup}

All transformer models were trained with HuggingFace Trainer: learning rate $2\times10^{-5}$,
batch size 32, linear schedule with 200 warmup steps, AdamW weight decay 0.01, up to 5 epochs
with early stopping (patience = 2) on validation weighted F1. Maximum sequence length: 64 tokens
(covers $>$99\% of tweets). Single random seed (42). All experiments ran on NVIDIA T4 GPU.

\subsection{Evaluation}

Primary metric: weighted F1 (accounts for class imbalance). Secondary: macro F1 (exposes
minority-class failure). Every system is evaluated on three subsets: full test, monolingual,
code-mixed.

\section{Experiments and Results}

\subsection{Main Results}

Table~\ref{tab:main} reports weighted F1 on the full test set.

\begin{table}[t]
\centering
\small
\setlength{\tabcolsep}{4pt}
\begin{tabular}{lccccc}
\toprule
\textbf{Model} & \textbf{hau} & \textbf{yor} & \textbf{ibo} & \textbf{pcm} & \textbf{swa} \\
\midrule
TF-IDF+LR & 69.84 & 72.64 & 77.35 & 65.06 & 57.83 \\
mBERT     & 74.25 & 66.05 & 75.78 & 64.33 & 52.67 \\
LAFT      & \textbf{76.34} & 70.27 & \textbf{77.24} & 68.02 & 60.86 \\
MAFT      & 75.57 & \textbf{70.27} & 75.02 & \textbf{69.00} & \textbf{62.51} \\
\bottomrule
\end{tabular}
\caption{Weighted F1 on full test sets. Bold = best per language.}
\label{tab:main}
\end{table}

\textbf{RQ1 (LAFT vs.\ MAFT).} LAFT achieves the highest weighted F1 on Hausa (76.34), the
highest-resource language (14,172 training examples, 1.07$\times$ imbalance). MAFT outperforms
LAFT on Swahili (62.51 vs.\ 60.86), the lowest-resource language (1,810 training examples).
This confirms the resource-driven crossover: monolingual specialisation is advantageous when
sufficient per-language data exists; cross-lingual parameter sharing compensates for data scarcity.

\textbf{mBERT vs.\ TF-IDF.} mBERT underperforms TF-IDF on Yor\`{u}b\'{a} (66.05 vs.\ 72.64).
Yor\`{u}b\'{a} is tonal; lexical features may capture sentiment-bearing diacritics more reliably
than subword tokenisation that fragments tonal characters. On Swahili, mBERT achieves only 52.67,
well below the TF-IDF baseline (57.83), highlighting the limitations of generic multilingual
pretraining for low-resource Bantu languages.

\textbf{RQ2 (class imbalance).} The weighted/macro F1 gap on Nigerian Pidgin is severe across
all models. MAFT achieves 69.00 weighted F1 but only 50.95 macro F1---a 18-point gap driven by
the 72-example neutral class being effectively ignored. Balanced class weights in TF-IDF and
standard cross-entropy in transformers both fail to overcome a 45$\times$ training imbalance.

\subsection{Code-Mixing Diagnostic}

Table~\ref{tab:cmi_results} reports the performance gap between monolingual and code-mixed subsets.

\begin{table}[t]
\centering
\small
\setlength{\tabcolsep}{3pt}
\begin{tabular}{llccc}
\toprule
\textbf{Lang.} & \textbf{Model} & \textbf{Mono.} & \textbf{Mixed} & \textbf{$\Delta$} \\
\midrule
hau & LAFT & 76.38 & 74.54 & $-1.84$ \\
hau & MAFT & 75.65 & 72.22 & $-3.43$ \\
yor & LAFT & 70.05 & 76.18 & $+6.13$ \\
yor & MAFT & 68.80 & 78.87 & $+10.07$ \\
ibo & LAFT & 76.93 & 81.48 & $+4.55$ \\
ibo & MAFT & 74.80 & 78.10 & $+3.30$ \\
pcm & LAFT & 65.96 & 71.25 & $+5.29$ \\
pcm & MAFT & 66.68 & 72.64 & $+5.96$ \\
swa & LAFT & 61.20 & 39.22 & $-21.98$\textsuperscript{*} \\
swa & MAFT & 62.61 & 55.39 & $-7.22$\textsuperscript{*} \\
\bottomrule
\end{tabular}
\caption{Weighted F1 on monolingual and code-mixed subsets.
$\Delta$ = Mixed $-$ Mono. \textsuperscript{*}Swahili mixed $n$=12.}
\label{tab:cmi_results}
\end{table}

\textbf{RQ3 (code-mixing effects).} Code-mixing does not universally degrade performance.
Yor\`{u}b\'{a} LAFT improves from 70.05 (mono) to 76.18 (mixed), a +6.13 gain. Igbo and
Nigerian Pidgin also show improvements on code-mixed subsets. A plausible explanation: English
function words in code-mixed Yor\`{u}b\'{a} and Igbo tweets carry strong sentiment polarity
that the model recognises, while purely African-language sentiment expressions involve tonal
and orthographic ambiguity that is harder to classify from romanised text.

Hausa shows modest degradation ($-$1.84 for LAFT), consistent with its near-monolingual
test distribution (2.4\% code-mixed). Swahili's severe LAFT degradation ($-$21.98) should
be treated with caution given only 12 mixed-subset tweets. MAFT's smaller Swahili degradation
($-$7.22) may indicate robustness from multilingual training, but cannot be confirmed from
12 examples.

\subsection{Error Analysis}

LIME explanations applied to TF-IDF misclassifications reveal four dominant error
categories. Code-switching errors account for 40\% of sampled Nigerian Pidgin
misclassifications (12/30), where the model encounters English phrases not associated with
surrounding Pidgin sentiment context. Orthographic variation (repeated characters,
non-standard spelling) explains errors in Hausa (2/30) and Igbo (2/30). Short
question-form tweets ($<$8 tokens with `?') contribute sarcasm-related errors across
all languages. The majority of remaining errors (labelled `other') likely reflect
label noise inherent in crowdsourced annotation of subjective sentiment.

\subsection{Responsible NLP Reflection}

Performance equity across languages is poor. Swahili's best weighted F1 (62.51) is
13.83 points below Hausa (76.34)---a gap driven entirely by data availability, not
linguistic complexity. The macro F1 gap on Nigerian Pidgin reveals that the neutral class
is invisible to all models, biasing any deployed system toward binary polarisation.

\section{Reflections and Discussion}

\textbf{What worked.} The LAFT/MAFT crossover is theoretically coherent and empirically
confirmed. TF-IDF's competitive performance on Yor\`{u}b\'{a} (72.64, above mBERT at 66.05)
suggests lexical baselines remain relevant for tonal languages, with practical implications
for GPU-constrained deployment.

\textbf{What did not work.} Nigerian Pidgin's neutral class was not rescued by any model.
The macro F1 ceiling near 50 indicates that 72 neutral training examples are insufficient
regardless of architecture. Data augmentation or targeted annotation is the most direct path
to improvement.

\textbf{Limitations.} Single random seed (42) without variance estimation limits
reproducibility confidence. The CMI filter is a heuristic with known noise on short tokens.
Swahili's 12-tweet mixed subset yields unreliable estimates. LIME was applied to TF-IDF
only; transformer interpretability was not evaluated. Hyperparameters were fixed without
grid search.

\textbf{Responsible AI.} All datasets are publicly available under Creative Commons
licences. Code and results are fully reproducible from the provided notebooks and
requirements files. AfriSenti was annotated by native speakers using majority-vote
protocols, though crowdsourced annotation inherits subjectivity biases. Future work should
extend to the full 14-language AfriSenti corpus and investigate data augmentation for
severely imbalanced classes.

\textbf{Extensions.} Three directions follow naturally. First, language-balanced MAFT
sampling (oversampling Swahili) may improve low-resource results while preserving
high-resource performance. Second, extending the CMI diagnostic to the full AfriSenti
corpus would provide sufficient sample sizes for robust Swahili analysis. Third, applying
LAFT/MAFT comparison at AfroXLMR-large scale would test whether the crossover persists at
greater model capacity.

\section{Conclusion}

We compared four sentiment systems across five AfriSenti languages using a novel
code-mixing diagnostic that partitions test sets into monolingual and code-mixed subsets.
Three findings emerge:

(1) LAFT outperforms MAFT on high-resource Hausa (76.34 vs.\ 75.57 wF1); MAFT outperforms
LAFT on low-resource Swahili (62.51 vs.\ 60.86), confirming a resource-driven crossover
(RQ1).

(2) Nigerian Pidgin's 45$\times$ class imbalance produces a severe wF1/mF1 gap that no
model resolves, demonstrating that data quality constraints dominate architectural choices
for severely imbalanced classes (RQ2).

(3) Code-mixing effects are language-dependent. Yor\`{u}b\'{a} and Igbo improve on
code-mixed subsets, while Hausa degrades modestly. The Swahili result is inconclusive
(RQ3).

For practitioners: LAFT is preferred for languages with $>$10,000 training examples; MAFT
offers measurable advantage below this threshold. The CMI diagnostic is a reusable
evaluation tool for any African language benchmark where English influence is present.

\bibliographystyle{plainnat}
\bibliography{references}

\end{document}
